\import{topology/topological-space.tex}
\import{set.tex}

\section{Separation}\label{form_sec_separation}

We call a space regular when points can be separated from closed sets; no
separation axiom is included in that adjective.  A \teethree-space is a
regular \teezero-space.  For \teeone we use the directional formulation;
quantification over ordered pairs supplies the reverse direction.

The unqualified separation adjectives below include the topological-space
domain. Their corresponding ``space'' phrases are natural-language aliases.

We begin with Kolmogorov, or \teezero, separation.

\begin{definition}\label{is_kolmogorov}
    $X$ is Kolmogorov iff $X$ is a topological space and
        for all $x,y\in\carrier[X]$ such that
            $x\neq y$ there exists $U\in\opens[X]$ such that
            $x\in U\not\ni y$ or $x\notin U\ni y$.
\end{definition}

\begin{classification}\label{kolmogorov_is_topological}
    Suppose $X$ is Kolmogorov.
    Then $X$ is a topological space.
\end{classification}
\begin{proof}
    Follows by \obv{is_kolmogorov}.
\end{proof}

\begin{abbreviation}\label{kolmogorov_space}
    $X$ is a Kolmogorov space iff $X$ is Kolmogorov.
\end{abbreviation}

\begin{classification}\label{kolmogorov_space_is_topological}
    Let $X$ be a Kolmogorov space.
    Then $X$ is a topological space.
\end{classification}
\begin{proof}
    Follows by \obv{is_kolmogorov}.
\end{proof}

\begin{abbreviation}\label{teezero}
    $X$ is \teezero\ iff $X$ is Kolmogorov.
\end{abbreviation}

\begin{abbreviation}\label{teezero_space}
    $X$ is a \teezero-space iff $X$ is a Kolmogorov space.
\end{abbreviation}

\begin{proposition}\label{kolmogorov_implies_kolmogorov_for_closeds}
    Suppose $X$ is a Kolmogorov space.
    Let $x,y\in\carrier[X]$.
    Suppose $x\neq y$.
    Then there exists $A\in\closeds{X}$ such that
        $x\in A\not\ni y$ or $x\notin A\ni y$.
\end{proposition}
\begin{proof}
    Take $U\in\opens[X]$ such that
        $x\in U\not\ni y$ or $x\notin U\ni y$
        by \obv{is_kolmogorov}.
    $\carrier[X]\setminus U\in\closeds{X}$
        by \obv{complement_of_open_elem_closeds}.
    $x\in(\carrier[X]\setminus U)\not\ni y$ or
        $x\notin(\carrier[X]\setminus U)\ni y$
        by \auto{setminus}.
\end{proof}

\begin{proposition}\label{kolmogorov_for_closeds_implies_kolmogorov}
    Let $X$ be a topological space.
    Suppose for all $x,y\in\carrier[X]$ such that $x\neq y$
        there exists $U\in\closeds{X}$ such that
        $x\in U\not\ni y$ or $x\notin U\ni y$.
    Then $X$ is Kolmogorov.
\end{proposition}
\begin{proof}
    Follows by \auto{closeds,is_kolmogorov,setminus}.
\end{proof}

\begin{proposition}\label{kolmogorov_iff_kolmogorov_for_closeds}
    Let $X$ be a topological space.
    $X$ is Kolmogorov iff for all $x,y\in\carrier[X]$ such that
        $x\neq y$ there exists $U\in\closeds{X}$ such that
        $x\in U\not\ni y$ or $x\notin U\ni y$.
\end{proposition}
\begin{proof}
    Follows by
        \auto{kolmogorov_implies_kolmogorov_for_closeds,kolmogorov_for_closeds_implies_kolmogorov}.
\end{proof}

The directional definition of \teeone separation is equivalent to the usual
two-neighbourhood formulation and to closed singletons.

\begin{definition}\label{teeone}
    $X$ is \teeone\ iff $X$ is a topological space and
        for all $x,y\in\carrier[X]$ such that
            $x\neq y$ there exists $U\in\opens[X]$ such that
            $x\in U\not\ni y$.
\end{definition}

\begin{classification}\label{teeone_is_topological}
    Suppose $X$ is \teeone.
    Then $X$ is a topological space.
\end{classification}
\begin{proof}
    Follows by \obv{teeone}.
\end{proof}

\begin{abbreviation}\label{teeone_space}
    $X$ is a \teeone-space iff $X$ is \teeone.
\end{abbreviation}

\begin{classification}\label{teeone_space_is_topological}
    Let $X$ be a \teeone-space.
    Then $X$ is a topological space.
\end{classification}
\begin{proof}
    Follows by \obv{teeone}.
\end{proof}

\begin{proposition}\label{teeone_iff_two_open_neighbourhoods}
    Let $X$ be a topological space.
    $X$ is \teeone\ iff for all $x,y\in\carrier[X]$ such that
        $x\neq y$ there exist $U,V\in\opens[X]$ such that
        $U\ni x\notin V$ and $V\ni y\notin U$.
\end{proposition}
\begin{proof}
    Show that if $X$ is \teeone, then for all
        $x,y\in\carrier[X]$ such that $x\neq y$
        there exist $U,V\in\opens[X]$ such that
        $U\ni x\notin V$ and $V\ni y\notin U$.
    \begin{subproof}
        Assume $X$ is \teeone.
        Fix $x,y\in\carrier[X]$.
        Assume $x\neq y$.
        Take $U\in\opens[X]$ such that $x\in U\not\ni y$
            by \obv{teeone}.
        Take $V\in\opens[X]$ such that $y\in V\not\ni x$
            by \obv{teeone}.
        Follows by \auto.
    \end{subproof}

    Show that if for all $x,y\in\carrier[X]$ such that $x\neq y$
        there exist $U,V\in\opens[X]$ such that
        $U\ni x\notin V$ and $V\ni y\notin U$, then $X$ is \teeone.
    \begin{subproof}
        Assume for all $x,y\in\carrier[X]$ such that $x\neq y$
            there exist $U,V\in\opens[X]$ such that
            $U\ni x\notin V$ and $V\ni y\notin U$.
        Follows by \auto{teeone}.
    \end{subproof}
    Follows by \auto.
\end{proof}

\begin{proposition}\label{teeone_implies_singletons_closed}
    Let $X$ be a \teeone-space.
    Suppose $x\in\carrier[X]$.
    Then $\{x\}$ is closed in $X$.
\end{proposition}
\begin{proof}
    $X$ is a topological space by notion.
    $\{x\}\subseteq\carrier[X]$ by notion.
    Show that $\carrier[X]\setminus\{x\}
        \subseteq\interior{\carrier[X]\setminus\{x\}}{X}$.
    \begin{subproof}
        Show that for every $y\in\carrier[X]\setminus\{x\}$ we have
            $y\in\interior{\carrier[X]\setminus\{x\}}{X}$.
        \begin{subproof}
            Fix $y\in\carrier[X]\setminus\{x\}$.
            Then $y\in\carrier[X]$ and $y\neq x$
                by \auto{setminus,singleton_iff}.
            Take $U\in\opens[X]$ such that $y\in U\not\ni x$
                by \obv{teeone}.
            $U\subseteq\carrier[X]\setminus\{x\}$
                by \auto{opens_type,subseteq,setminus,singleton_iff}.
            Follows by \obv{interior_elem_intro}.
        \end{subproof}
        Follows by \obv{subseteq}.
    \end{subproof}
    $\interior{\carrier[X]\setminus\{x\}}{X}
        \subseteq\carrier[X]\setminus\{x\}$
        by \obv{interior_subseteq}.
    Hence $\carrier[X]\setminus\{x\}
        =\interior{\carrier[X]\setminus\{x\}}{X}$
        by \obv{subseteq_antisymmetric}.
    $\interior{\carrier[X]\setminus\{x\}}{X}$ is open in $X$
        by \obv{is_open,interior_is_open}.
    Thus $\carrier[X]\setminus\{x\}$ is open in $X$.
    Hence $\carrier[X]\setminus\{x\}\in\opens[X]$ by \obv{is_open}.
    $\{x\}\in\pow{\carrier[X]}$ by \obv{powerset_intro}.
    Hence $\{x\}\in\closeds{X}$ by \obv{closeds}.
    Follows by \obv{is_closed_in}.
\end{proof}

\begin{proposition}\label{singletons_closed_implies_teeone}
    Let $X$ be a topological space.
    Suppose for every $x\in\carrier[X]$ we have
        $\{x\}$ is closed in $X$.
    Then $X$ is \teeone.
\end{proposition}
\begin{proof}
    Show that for all $x,y\in\carrier[X]$ such that $x\neq y$
        there exists $U\in\opens[X]$ such that $x\in U\not\ni y$.
    \begin{subproof}
        Fix $x,y\in\carrier[X]$.
        Assume $x\neq y$.
        $\{y\}$ is closed in $X$.
        Hence $\{y\}\in\closeds{X}$ by \obv{is_closed_in}.
        Then $\carrier[X]\setminus\{y\}\in\opens[X]$ by \obv{closeds}.
        Hence $\carrier[X]\setminus\{y\}$ is open in $X$ by \obv{is_open}.
        We have $x\in\carrier[X]\setminus\{y\}$ and
            $y\notin\carrier[X]\setminus\{y\}$
            by \auto{setminus,singleton_iff}.
        Follows by \obv{is_open}.
    \end{subproof}
    Follows by \obv{teeone}.
\end{proof}

\begin{theorem}\label{teeone_iff_singletons_closed}
    Let $X$ be a topological space.
    $X$ is \teeone\ iff for every $x\in\carrier[X]$ we have
        $\{x\}$ is closed in $X$.
\end{theorem}
\begin{proof}
    Follows by \auto{teeone_implies_singletons_closed,singletons_closed_implies_teeone}.
\end{proof}

Hausdorff, or \teetwo, separation requires disjoint neighbourhoods.

\begin{definition}\label{is_hausdorff}
    $X$ is Hausdorff iff $X$ is a topological space and
        for all $x,y\in\carrier[X]$ such that
            $x\neq y$ there exist $U,V\in\opens[X]$ such that
            $x\in U$ and $y\in V$ and $U$ is disjoint from $V$.
\end{definition}

\begin{classification}\label{hausdorff_is_topological}
    Suppose $X$ is Hausdorff.
    Then $X$ is a topological space.
\end{classification}
\begin{proof}
    Follows by \obv{is_hausdorff}.
\end{proof}

\begin{abbreviation}\label{hausdorff_space}
    $X$ is a Hausdorff space iff $X$ is Hausdorff.
\end{abbreviation}

\begin{classification}\label{hausdorff_space_is_topological}
    Let $X$ be a Hausdorff space.
    Then $X$ is a topological space.
\end{classification}
\begin{proof}
    Follows by \obv{is_hausdorff}.
\end{proof}

\begin{abbreviation}\label{teetwo}
    $X$ is \teetwo\ iff $X$ is Hausdorff.
\end{abbreviation}

\begin{abbreviation}\label{teetwo_space}
    $X$ is a \teetwo-space iff $X$ is a Hausdorff space.
\end{abbreviation}

\begin{proposition}\label{teeone_space_is_teezero_space}
    Let $X$ be a \teeone-space.
    Then $X$ is a \teezero-space.
\end{proposition}
\begin{proof}
    Follows by \auto{is_kolmogorov,teeone}.
\end{proof}

\begin{proposition}\label{teetwo_space_is_teeone_space}
    Let $X$ be a \teetwo-space.
    Then $X$ is a \teeone-space.
\end{proposition}
\begin{proof}
    Follows by \auto{is_hausdorff,teeone,disjoint}.
\end{proof}

\begin{definition}\label{is_regular}
    $X$ is regular iff $X$ is a topological space and
        for all $C,p$ such that
            $p\in\carrier[X]$ and $p\notin C\in\closeds{X}$ there exist
            $U,V\in\opens[X]$ such that $p\in U$ and $C\subseteq V$ and
            $U\inter V=\emptyset$.
\end{definition}

\begin{classification}\label{regular_is_topological}
    Suppose $X$ is regular.
    Then $X$ is a topological space.
\end{classification}
\begin{proof}
    Follows by \obv{is_regular}.
\end{proof}

\begin{abbreviation}\label{regular_space}
    $X$ is a regular space iff $X$ is regular.
\end{abbreviation}

\begin{classification}\label{regular_space_is_topological}
    Let $X$ be a regular space.
    Then $X$ is a topological space.
\end{classification}
\begin{proof}
    Follows by \obv{is_regular}.
\end{proof}


\begin{abbreviation}\label{teethree}
    $X$ is \teethree\ iff $X$ is regular and $X$ is \teezero.
\end{abbreviation}

\begin{abbreviation}\label{teethree_space}
    $X$ is a \teethree-space iff $X$ is \teethree.
\end{abbreviation}

\begin{classification}\label{teethree_space_is_topological}
    Let $X$ be a \teethree-space.
    Then $X$ is a topological space.
\end{classification}
\begin{proof}
    Follows by \obv{is_regular}.
\end{proof}

\begin{proposition}\label{regular_implies_closed_neighbourhood_in_open}
    Let $X$ be a regular space.
    For every $U\in\opens[X]$ we have for every $x\in U$ there exists
        $N\in\neighbourhoods{x}{X}$ such that
        $N\subseteq U$ and $N$ is closed in $X$.
\end{proposition}
\begin{proof}
    $X$ is a topological space by \obv{is_regular}.
    Fix $U\in\opens[X]$.
    Fix $x\in U$.
    $U\subseteq\carrier[X]$ by \obv{opens_type}.
    Let $C=\carrier[X]\setminus U$.
    $C\in\closeds{X}$ by \auto{complement_of_open_elem_closeds}.
    $x\in\carrier[X]$ and $x\notin C$ by \auto{subseteq,setminus}.
    Take $A,B\in\opens[X]$ such that $x\in A$ and $C\subseteq B$ and
        $A\inter B=\emptyset$ by \obv{is_regular}.
    $A\subseteq\carrier[X]$ and $B\subseteq\carrier[X]$
        by \obv{opens_type}.
    $A\subseteq\carrier[X]\setminus B$
        by \obv{subseteq_setminus}.
    Let $N=\carrier[X]\setminus B$.
    $N\subseteq\carrier[X]$ by \auto{setminus_subseteq}.
    Hence $N\in\pow{\carrier[X]}$ by \obv{powerset_intro}.
    $B$ is open in $X$ by \obv{is_open}.
    $N$ is closed in $X$ by \auto{complement_of_open_is_closed}.
    $x\in A\subseteq N$.
    Hence $N\in\neighbourhoods{x}{X}$ by \auto{neighbourhoods}.
    $\carrier[X]\setminus B\subseteq\carrier[X]\setminus C$
        by \obv{subseteq_implies_setminus_supseteq}.
    $\carrier[X]\setminus C=U$ by \auto{double_relative_complement}.
    Hence $N\subseteq U$.
\end{proof}

\begin{theorem}\label{regular_iff_closed_neighbourhood_in_open}
    Let $X$ be a topological space.
    $X$ is regular iff for every $U\in\opens[X]$ we have for every $x\in U$
        there exists $N\in\neighbourhoods{x}{X}$ such that
        $N\subseteq U$ and $N$ is closed in $X$.
\end{theorem}
\begin{proof}
    Show that if $X$ is regular, then for every $U\in\opens[X]$ we have
        for every $x\in U$ there exists $N\in\neighbourhoods{x}{X}$ such that
        $N\subseteq U$ and $N$ is closed in $X$.
    \begin{subproof}
        Assume $X$ is regular.
        Follows by \obv{regular_implies_closed_neighbourhood_in_open}.
    \end{subproof}

    Show that if for every $U\in\opens[X]$ we have for every $x\in U$
        there exists $N\in\neighbourhoods{x}{X}$ such that
        $N\subseteq U$ and $N$ is closed in $X$, then $X$ is regular.
    \begin{subproof}
        Assume for every $U\in\opens[X]$ we have for every $x\in U$
            there exists $N\in\neighbourhoods{x}{X}$ such that
            $N\subseteq U$ and $N$ is closed in $X$.
        Show that for all $C,p$ such that $p\in\carrier[X]$ and
            $p\notin C\in\closeds{X}$ there exist $A,V\in\opens[X]$
            such that $p\in A$ and $C\subseteq V$ and
            $A\inter V=\emptyset$.
        \begin{subproof}
            Fix $C,p$.
            Assume $p\in\carrier[X]$ and $p\notin C\in\closeds{X}$.
            $C\in\pow{\carrier[X]}$ by \obv{closeds}.
            Hence $C\subseteq\carrier[X]$ by \obv{powerset_elim}.
            $C$ is closed in $X$ by \obv{is_closed_in}.
            Let $U=\carrier[X]\setminus C$.
            $U$ is open in $X$ by \obv{complement_of_closed_is_open}.
            Hence $U\in\opens[X]$ by \obv{is_open}.
            $p\in U$ by \obv{setminus}.
            Take $N\in\neighbourhoods{p}{X}$ such that $N\subseteq U$
                and $N$ is closed in $X$.
            $N\in\pow{\carrier[X]}$ by \obv{neighbourhoods}.
            Hence $N\subseteq\carrier[X]$ by \obv{powerset_elim}.
            Take $A\in\opens[X]$ such that $p\in A\subseteq N$
                by \obv{neighbourhoods}.
            Let $V=\carrier[X]\setminus N$.
            $V$ is open in $X$ by \auto{complement_of_closed_is_open}.
            Hence $V\in\opens[X]$ by \obv{is_open}.
            $C\inter N=\emptyset$
                by \auto{inter,setext,emptyset,setminus,subseteq}.
            Hence $C\subseteq V$ by \obv{subseteq_setminus}.
            $A\inter V=\emptyset$
                by \auto{inter,setext,emptyset,setminus,subseteq}.
            Follows by \auto.
        \end{subproof}
        Follows by \obv{is_regular}.
    \end{subproof}
    Follows by \auto.
\end{proof}

\begin{corollary}\label{teethree_implies_closed_neighbourhood_in_open}
    Let $X$ be a \teethree-space.
    For every $U\in\opens[X]$ we have for every $x\in U$ there exists
        $N\in\neighbourhoods{x}{X}$ such that
        $N\subseteq U$ and $N$ is closed in $X$.
\end{corollary}
\begin{proof}
    Follows by \obv{regular_implies_closed_neighbourhood_in_open}.
\end{proof}

\begin{definition}[Closed neighbourhoods]\label{closed_neighbourhoods}
    $\closedNeighbourhoods{H}{X}
        =\{N\in\neighbourhoodsSet{H}{X}\mid N\in\closeds{X}\}$.
\end{definition}

\begin{proposition}\label{carrier_is_closed_neighbourhood}
    Let $X$ be a topological space.
    Suppose $H\subseteq\carrier[X]$.
    Then $\carrier[X]\in\closedNeighbourhoods{H}{X}$.
\end{proposition}
\begin{proof}
    $\carrier[X]\in\pow{\carrier[X]}$ by \obv{powerset_top}.
    $\carrier[X]\in\opens[X]$ by \obv{carrier_open}.
    Hence $\carrier[X]\in\neighbourhoodsSet{H}{X}$
        by \obv{neighbourhoods_set,subseteq_refl}.
    $\carrier[X]$ is closed in $X$ by \obv{carrier_is_closed}.
    Hence $\carrier[X]\in\closeds{X}$ by \obv{is_closed_in}.
    Follows by \obv{closed_neighbourhoods}.
\end{proof}

\begin{theorem}\label{regular_iff_closed_is_intersection_of_closed_neighbourhoods}
    Let $X$ be a topological space.
    $X$ is regular iff for every $H\in\closeds{X}$ we have
        $H=\inters{\closedNeighbourhoods{H}{X}}$.
\end{theorem}
\begin{proof}
    Show that if $X$ is regular, then for every $H\in\closeds{X}$ we have
        $H=\inters{\closedNeighbourhoods{H}{X}}$.
    \begin{subproof}
        Assume $X$ is regular.
        Fix $H\in\closeds{X}$.
        $H\in\pow{\carrier[X]}$ by \obv{closeds}.
        Hence $H\subseteq\carrier[X]$ by \obv{powerset_elim}.
        $\carrier[X]\in\closedNeighbourhoods{H}{X}$
            by \obv{carrier_is_closed_neighbourhood}.
        Hence $\closedNeighbourhoods{H}{X}$ is inhabited.

        Show that $H\subseteq\inters{\closedNeighbourhoods{H}{X}}$.
        \begin{subproof}
            Show that for every $N\in\closedNeighbourhoods{H}{X}$
                we have $H\subseteq N$.
            \begin{subproof}
                Fix $N\in\closedNeighbourhoods{H}{X}$.
                Then $N\in\neighbourhoodsSet{H}{X}$
                    by \obv{closed_neighbourhoods}.
                Take $V\in\opens[X]$ such that $H\subseteq V\subseteq N$
                    by \obv{neighbourhoods_set}.
                Follows by \obv{subseteq_transitive}.
            \end{subproof}
            Follows by \obv{inters_greatest}.
        \end{subproof}

        Show that $\inters{\closedNeighbourhoods{H}{X}}\subseteq H$.
        \begin{subproof}
            Show that for every
                $x\in\inters{\closedNeighbourhoods{H}{X}}$ we have $x\in H$.
            \begin{subproof}
                Fix $x\in\inters{\closedNeighbourhoods{H}{X}}$.
                Then $x\in\carrier[X]$
                    by \obv{inters_destr,carrier_is_closed_neighbourhood}.
                Suppose not.
                Take $A,B\in\opens[X]$ such that $x\in A$ and $H\subseteq B$
                    and $A\inter B=\emptyset$ by \obv{is_regular,closeds}.
                $A\subseteq\carrier[X]$ and $B\subseteq\carrier[X]$
                    by \obv{opens_type}.
                $B\inter A=\emptyset$ by \obv{inter_comm}.
                Hence $B\subseteq\carrier[X]\setminus A$
                    by \obv{subseteq_setminus}.
                Let $M=\carrier[X]\setminus A$.
                $M\in\closeds{X}$ by \auto{complement_of_open_elem_closeds}.
                $M\in\pow{\carrier[X]}$ by \obv{closeds}.
                Hence $M\in\neighbourhoodsSet{H}{X}$
                    by \auto{neighbourhoods_set}.
                Thus $M\in\closedNeighbourhoods{H}{X}$
                    by \obv{closed_neighbourhoods}.
                Hence $x\in M$ by \obv{inters_destr}.
                $x\notin M$ by \auto{setminus}.
                Contradiction.
            \end{subproof}
            Follows by \obv{subseteq}.
        \end{subproof}
        Follows by \obv{subseteq_antisymmetric}.
    \end{subproof}

    Show that if for every $H\in\closeds{X}$ we have
        $H=\inters{\closedNeighbourhoods{H}{X}}$, then $X$ is regular.
    \begin{subproof}
        Assume for every $H\in\closeds{X}$ we have
            $H=\inters{\closedNeighbourhoods{H}{X}}$.
        Show that for all $C,p$ such that $p\in\carrier[X]$ and
            $p\notin C\in\closeds{X}$ there exist $U,V\in\opens[X]$
            such that $p\in U$ and $C\subseteq V$ and
            $U\inter V=\emptyset$.
        \begin{subproof}
            Fix $C,p$.
            Assume $p\in\carrier[X]$ and $p\notin C\in\closeds{X}$.
            $C\in\pow{\carrier[X]}$ by \obv{closeds}.
            Hence $C\subseteq\carrier[X]$ by \obv{powerset_elim}.
            $\carrier[X]\in\closedNeighbourhoods{C}{X}$
                by \obv{carrier_is_closed_neighbourhood}.
            Hence $\closedNeighbourhoods{C}{X}$ is inhabited.
            $p\notin\inters{\closedNeighbourhoods{C}{X}}$.
            Take $N\in\closedNeighbourhoods{C}{X}$ such that $p\notin N$
                by \auto{inters_iff_forall}.
            $N\in\neighbourhoodsSet{C}{X}$ and $N\in\closeds{X}$
                by \obv{closed_neighbourhoods}.
            $N\in\pow{\carrier[X]}$ by \obv{closeds}.
            $N$ is closed in $X$ by \obv{is_closed_in}.
            Hence $N\subseteq\carrier[X]$ by \obv{powerset_elim}.
            Take $V\in\opens[X]$ such that $C\subseteq V\subseteq N$
                by \obv{neighbourhoods_set}.
            Let $U=\carrier[X]\setminus N$.
            $U$ is open in $X$ by \auto{complement_of_closed_is_open}.
            Hence $U\in\opens[X]$ by \obv{is_open}.
            $p\in U$ by \obv{setminus}.
            $U\inter V=\emptyset$
                by \auto{inter,setext,emptyset,setminus,subseteq}.
            Follows by \auto.
        \end{subproof}
        Follows by \obv{is_regular}.
    \end{subproof}
    Follows by \auto.
\end{proof}

\begin{corollary}\label{teethree_implies_closed_is_intersection_of_closed_neighbourhoods}
    Let $X$ be a \teethree-space.
    For every $H\in\closeds{X}$ we have
        $H=\inters{\closedNeighbourhoods{H}{X}}$.
\end{corollary}
\begin{proof}
    $X$ is regular by definition.
    Follows by
        \obv{is_regular,
            regular_iff_closed_is_intersection_of_closed_neighbourhoods}.
\end{proof}


\begin{theorem}\label{teethree_space_is_teetwo_space}
    Let $X$ be a \teethree-space.
    Then $X$ is a \teetwo-space.
\end{theorem}
\begin{proof}
    Show that $X$ is Hausdorff.
    \begin{subproof}
        Show that for all $x,y\in\carrier[X]$ such that $x\neq y$
            there exist $U,V\in\opens[X]$ such that $x\in U$ and
            $y\in V$ and $U$ is disjoint from $V$.
        \begin{subproof}
            Fix $x,y\in\carrier[X]$.
            Assume $x\neq y$.
            Take $C\in\closeds{X}$ such that
                $x\in C\not\ni y$ or $x\notin C\ni y$
                by \obv{kolmogorov_implies_kolmogorov_for_closeds}.
            \begin{byCase}
            \caseOf{$x\in C$.}
                Hence $y\notin C$.
                Take $V,U\in\opens[X]$ such that $y\in V$ and
                    $C\subseteq U$ and $V\inter U=\emptyset$
                    by \obv{is_regular}.
                Then $x\in U$ by \obv{subseteq}.
                $V$ is disjoint from $U$
                    by \auto{disjoint,inter,emptyset}.
                Hence $U$ is disjoint from $V$
                    by \obv{disjoint_symmetric}.
                Follows by \auto.
            \caseOf{$x\notin C$.}
                Hence $y\in C$.
                Take $U,V\in\opens[X]$ such that $x\in U$ and
                    $C\subseteq V$ and $U\inter V=\emptyset$
                    by \obv{is_regular}.
                Then $y\in V$ by \obv{subseteq}.
                $U$ is disjoint from $V$
                    by \auto{disjoint,inter,emptyset}.
                Follows by \auto.
            \end{byCase}
        \end{subproof}
        Follows by \obv{is_hausdorff}.
    \end{subproof}
\end{proof}

\begin{proposition}\label{indiscrete_space_is_regular}
    Let $X$ be a topological space.
    Suppose $\opens[X]=\{\emptyset,\carrier[X]\}$.
    Then $X$ is regular.
\end{proposition}
\begin{proof}
    Show that for all $C,p$ such that $p\in\carrier[X]$ and
        $p\notin C\in\closeds{X}$ there exist $U,V\in\opens[X]$
        such that $p\in U$ and $C\subseteq V$ and
        $U\inter V=\emptyset$.
    \begin{subproof}
        Fix $C,p$.
        Assume $p\in\carrier[X]$ and $p\notin C\in\closeds{X}$.
        $C\in\pow{\carrier[X]}$ by \obv{closeds}.
        $C$ is closed in $X$ by \obv{is_closed_in}.
        Hence $C\subseteq\carrier[X]$ by \obv{powerset_elim}.
        $\carrier[X]\setminus C$ is open in $X$
            by \obv{complement_of_closed_is_open}.
        Hence $\carrier[X]\setminus C\in\opens[X]$ by \obv{is_open}.
        Hence $\carrier[X]\setminus C=\emptyset$ or
            $\carrier[X]\setminus C=\carrier[X]$
            by \auto{cons_iff,singleton_iff}.
        $p\in\carrier[X]\setminus C$ by \obv{setminus}.
        Hence $\carrier[X]\setminus C\neq\emptyset$ by \obv{emptyset}.
        Thus $\carrier[X]\setminus C=\carrier[X]$.
        $\carrier[X]\setminus(\carrier[X]\setminus C)=C$
            by \obv{double_relative_complement}.
        Hence $C=\emptyset$ by \auto{setminus_self}.
        $\carrier[X],\emptyset\in\opens[X]$
            by \auto{carrier_open,emptyset_open}.
        $\carrier[X]\inter\emptyset=\emptyset$ by \obv{inter_emptyset}.
        Follows by \auto{emptyset_subseteq}.
    \end{subproof}
    Follows by \obv{is_regular}.
\end{proof}

\begin{proposition}\label{nontrivial_indiscrete_space_is_not_kolmogorov}
    Let $X$ be a topological space.
    Suppose $\opens[X]=\{\emptyset,\carrier[X]\}$.
    Suppose $a,b\in\carrier[X]$ and $a\neq b$.
    Then $X$ is not Kolmogorov.
\end{proposition}
\begin{proof}
    Suppose $X$ is Kolmogorov.
    Take $U\in\opens[X]$ such that
        $a\in U\not\ni b$ or $a\notin U\ni b$
        by \obv{is_kolmogorov}.
    Then $U=\emptyset$ or $U=\carrier[X]$
        by \auto{cons_iff,singleton_iff}.
    \begin{byCase}
    \caseOf{$U=\emptyset$.}
        Then $a,b\notin U$ by \obv{emptyset}.
        Contradiction.
    \caseOf{$U\neq\emptyset$.}
        Hence $U=\carrier[X]$.
        Thus $a,b\in U$.
        Contradiction.
    \end{byCase}
\end{proof}

The preceding statements have concrete content: the canonical two-point
indiscrete space is a checked exact record, rather than an assumed abstract
space.

\begin{proposition}\label{two_point_indiscrete_space_is_regular}
    $\TwoPointIndiscreteSpace$ is regular.
\end{proposition}
\begin{proof}
    $\carrier[\TwoPointIndiscreteSpace]=\twoPointCarrier$ by
        \obv{two_point_indiscrete_spaceprojection_carrier}.
    $\opens[\TwoPointIndiscreteSpace]
        =\indiscreteOpens{\twoPointCarrier}$ by
        \obv{two_point_indiscrete_spaceprojection_opens}.
    Hence $\opens[\TwoPointIndiscreteSpace]
        =\{\emptyset,\carrier[\TwoPointIndiscreteSpace]\}$
        by \auto{indiscrete_opens}.
    Follows by \obv{two_point_indiscrete_space,indiscrete_space_is_regular}.
\end{proof}

\begin{proposition}\label{two_point_indiscrete_space_is_not_kolmogorov}
    $\TwoPointIndiscreteSpace$ is not Kolmogorov.
\end{proposition}
\begin{proof}
    $\carrier[\TwoPointIndiscreteSpace]=\twoPointCarrier$ by
        \obv{two_point_indiscrete_spaceprojection_carrier}.
    $\opens[\TwoPointIndiscreteSpace]
        =\indiscreteOpens{\twoPointCarrier}$ by
        \obv{two_point_indiscrete_spaceprojection_opens}.
    Hence $\opens[\TwoPointIndiscreteSpace]
        =\{\emptyset,\carrier[\TwoPointIndiscreteSpace]\}$
        by \auto{indiscrete_opens}.
    $\emptyset,\{\emptyset\}\in\carrier[\TwoPointIndiscreteSpace]$
        by \auto{two_point_carrier,cons_iff,singleton_iff}.
    $\emptyset\in\{\emptyset\}$ by \auto{singleton_iff}.
    $\emptyset\notin\emptyset$ by \obv{emptyset}.
    Hence $\emptyset\neq\{\emptyset\}$.
    Follows by
        \obv{two_point_indiscrete_space,nontrivial_indiscrete_space_is_not_kolmogorov}.
\end{proof}
